% Lösungen Einheit 2 von 4 – Terme zusammenfassen (zusammenbau v0.9)
\einheitenkopf{Einheit 2 von 4 · Terme zusammenfassen}

\erg{15}{a) $4y$ und $9y$ \quad b) $2a$ und $a$; $6$ und $8$ \quad c) $5x^2$ und $3x^2$ \quad d) $7b$ und $5b$}
\erg{16}{a) $1$ (denn $a = 1a$) \quad b) $-1$ (denn $-b = -1b$) \quad c) $1{,}5$ \quad d) $-4$}
\erg{17}{a) $+5a$, $-2b$, $+7$ \quad b) $-3x$, $+6$, $-y$ \quad c) $+2m$, $-9$, $-5k$ \quad d) $-6$, $+3c$, $-4d$}
\erg{18}{a) Vorzahlen addieren: $7x + 2x = 9x$ \quad b) Vorzahlen addieren: $7x + 4x = 11x$ \quad c) Vorzahlen addieren: $7x + 5x = 12x$ \quad d) Vorzahlen addieren: $7x + 8x = 15x$ \quad e) $5x$ \quad f) $9x$ \quad g) $7x + 3$}
\erg{19}{a) $6a + 3b$ \quad b) $5y$ \quad c) $-5x$ \quad d) $5x$ \quad e) $6x^2 + 5x$ \quad f) $5x$}
\erg{20}{$8a - 5a = 3a$, Term $3a - 2a^2$; für $a = 3$: $3 \cdot 3 - 2 \cdot 3^2 = 9 - 18 = -9$}
\erg{21}{a) Zahlen malnehmen: $5 \cdot 2x = 10x$ \quad b) Zahlen malnehmen: $5 \cdot 3x = 15x$ \quad c) Zahlen malnehmen: $5 \cdot 6x = 30x$ \quad d) Zahlen malnehmen: $5 \cdot 7x = 35x$ \quad e) $20x$ \quad f) $8y^2$ \quad g) $-20x$ \quad h) $20ab$ \quad i) $30x$}
\erg{22}{$-12ab$}
\erg{23}{a) Die Zahl 3 ist kein $a$-Glied und wurde mitgezählt. Richtig: $5a + 3 + 2a = 7a + 3$. \quad b) Richtig. $b$ hat die Vorzahl 1, also $1b + 8b = 9b$; die Zahl $-3$ ist kein $b$-Glied und bleibt stehen. \quad c) Nicht stimmen können $9x^2$, $5y$ und $6b$. $9x^2$: Beim Zusammenfassen bleibt die Hochzahl gleich. $5y$: 6 minus 11 ist negativ, die Vorzahl muss negativ sein. $6b$: $b$ mal $b$ ergibt $b^2$, die Hochzahl muss 2 sein.}
\erg{24}{a) Nein: $x$ und $y$ sind verschiedene Variablen, gleiche Vorzahlen reichen nicht. \quad b) (1) wahr, denn man addiert nur die Vorzahlen, die Variable bleibt. (2) falsch, z. B. $5x + 5y$: verschiedene Variablen. (3) wahr, denn beim Addieren bleibt die Variable mit ihrer Hochzahl gleich, z. B. $2x^2 + 4x^2 = 6x^2$. \quad c) Nein; man rechnet $3 - 8$, nicht $8 - 3$: Wer mehr abzieht, als da ist, landet unter null, also $3y - 8y = -5y$.}
\erg{25}{a) $3x + 2 + 3x + 2 = 6x + 4$ \quad b) $4a + 3a + 2a = 9a$ \quad c) Zum Beispiel ein gleichseitiges Dreieck mit der Seitenlänge $3a$ cm. Zusammenfassen: $3a + 3a + 3a = 9a$.}
\erg{26}{a) Ja; Term aufstellen: $45p + 38p + 52p$; zusammenfassen: $135p$; einsetzen: $135 \cdot 0{,}40 = 54$ €, das sind 4 € mehr als 50 €. \quad b) $s + 2s + s + 2s = 6s$; für $s = 8$: $6 \cdot 8 = 48$ m \quad c) $85n + 24n + 36n = 145n$; für $n = 26$: $145 \cdot 26 = 3\,770$ €}
